% 中文全要素：摘要、关键词、参考文献、定理环境都必须是中文措辞
\documentclass[UTF8]{ctexart}
\usepackage{amsmath,amsthm}

\newtheorem{theorem}{定理}
\newtheorem{lemma}[theorem]{引理}

\begin{document}

\title{含完整中文结构的稿件}
\author{张三 \and 李四}
\date{}

\begin{abstract}
本文检查中文稿件里摘要、关键词与参考文献的措辞是否跟随稿件语言，
而不是一律印成英文标题。
\end{abstract}

\keywords{拉泰赫转Word，公式对象，中文排版}

\section{引言}
\label{sec:cntro}

这是一个中文段落，用来验证正文排版。行内公式 $a+b$ 与行间公式如下：

\begin{equation}
\int_{0}^{1} f(x)\,\mathrm{d}x = F(1) - F(0)
\label{eq:cnint}
\end{equation}

\begin{theorem}
\label{thm:cn}
中文定理环境应当显示中文名。
\end{theorem}

\begin{lemma}
\label{lem:cn}
与定理共用计数器的引理。
\end{lemma}

见第~\ref{sec:cntro} 节、式~\eqref{eq:cnint} 与定理~\ref{thm:cn}。

\begin{thebibliography}{9}
\bibitem{zk1} 王五. 中文文献条目一. 2020.
\bibitem{zk2} 赵六. 中文文献条目二. 2021.
\end{thebibliography}

\end{document}
